\documentclass[11pt]{article}
\usepackage[margin=1in]{geometry}
\usepackage{amsmath,amssymb,amsthm}
\usepackage{xurl}
\usepackage{hyperref}
\sloppy
\let\oldus\_ \renewcommand{\_}{\oldus\allowbreak}

\newtheorem{theorem}{Theorem}
\newtheorem{lemma}[theorem]{Lemma}
\theoremstyle{remark}
\newtheorem{remark}[theorem]{Remark}
\theoremstyle{definition}
\newtheorem{definition}[theorem]{Definition}

\title{Conjecture 00000008228 is false:\\ the Ehrhart function of a rank-one weight polytope is never monic}
\author{Jackmeson1}
\date{}

\begin{document}
\maketitle

\section{The conjecture}

The English text of conjecture 00000008228 reads:

\begin{quote}
Definition: The extremities of weight vectors in geometric Satake: the vertices sat\_vertex of the
weight polytope under the Satake equivalence on the affine Grassmannian. Conjecture: sat\_vertex (the
vertices of the Satake polytope) are the weights of Mirkovic-Vilonen cycles (an MV-weight law); the
volume of the polytope is the normalization of the representation dimension over $|W|$ (a
volume-dimension normalization law); the count of extremal points (vertex count) equals the count of
extremal elements of the saturated set of the weight lattice (an extremal count law); and the
Ehrhart function of the polytope is a monic polynomial of degree the semisimple rank (a semisimple
rank law).
\end{quote}

The conjecture is a conjunction of four laws. We refute the fourth one, the \emph{semisimple rank
law}: ``the Ehrhart function of the polytope is a monic polynomial of degree the semisimple rank''.
Hence the conjunction is false. We do not use, and make no claim about, the first three laws
(Remark~\ref{rem:clause2} comments on the second one but is not part of the proof).

\section{Reading and definitions}

No group is fixed by the text, so the laws are asserted for a reductive group $G$ and a coweight
$\lambda$ in general; one pair $(G,\lambda)$ violating the fourth law refutes it. We take $G = SL_2$
and $G = PGL_2$.

\begin{definition}[Root datum, Weyl group]
A root datum is a pair of free abelian groups $X, X^\vee$ of finite rank in perfect pairing
$\langle\cdot,\cdot\rangle$, with roots $\Phi\subset X$ and coroots $\Phi^\vee\subset X^\vee$,
$\langle\alpha,\alpha^\vee\rangle = 2$, and both sets stable under the reflections
$s_\alpha(x) = x - \langle x,\alpha^\vee\rangle\alpha$ (and dually). The Weyl group $W$ is the group
generated by the reflections. In Lean this is Mathlib's \texttt{RootDatum} (a \texttt{RootPairing}
over $\mathbb{Z}$) and \texttt{RootPairing.weylGroup}, the subgroup generated by the reflections.
\end{definition}

For $a,b\in\mathbb{Z}$ with $ab = 2$, let $\mathcal{R}(a,b)$ be the root datum with
$X = X^\vee = \mathbb{Z}$, pairing $\langle x,y\rangle = xy$, roots $\pm a$ and coroots $\pm b$
(Lean: \texttt{rk1 a b hab}). The datum $(a,b) = (2,1)$ is that of $SL_2$
(\texttt{rootDatumSL2}) and $(a,b) = (1,2)$ is that of $PGL_2$ (\texttt{rootDatumPGL2}). In geometric
Satake the polytope lives in the cocharacter lattice of $G$, which is the character lattice of the
dual group $G^\vee$. For $G = PGL_2$ this is the $SL_2$ datum, and for $G = SL_2$ it is the $PGL_2$
datum. All four solutions of $ab = 2$ are covered, so this choice does not matter.

\begin{definition}[Weight polytope]
For $\lambda\in X$, the weight polytope is
$P_\lambda = \operatorname{conv}(W\cdot\lambda)\subset X\otimes\mathbb{R}$. For $X = \mathbb{Z}$ we
identify $X\otimes\mathbb{R} = \mathbb{R}$ (Lean: \texttt{weightPolytope P lam}, the real convex hull of
the image of \texttt{MulAction.orbit P.weylGroup lam} under $\mathbb{Z}\to\mathbb{R}$).
\end{definition}

\begin{definition}[Ehrhart function]
Following \cite{wiki}: ``Informally, if P is a polytope, and tP is the polytope formed by expanding P
by a factor of t in each dimension, then L(P, t) is the number of integer lattice points in tP.''
For a lattice $\Lambda\le X$ and $t\in\mathbb{Z}_{>0}$ we set
$L_\Lambda(P,t) = \#\{x\in\Lambda : x\in tP\}$ (Lean: \texttt{ehrhart Lam Q t}, using
\texttt{Set.ncard}; all sets counted here are finite).
\end{definition}

\begin{definition}[Semisimple rank]
The semisimple rank is the rank of the root lattice $\mathbb{Z}\Phi\le X$ (Lean:
\texttt{semisimpleRank P}, the \texttt{Module.finrank} over $\mathbb{Z}$ of
\texttt{P.rootSpan}~$\mathbb{Z}$). For $\mathcal{R}(a,b)$ it equals $1$ (\texttt{semisimpleRank\_rk1}).
\end{definition}

\paragraph{Conventions.}
(i) Which lattice the Ehrhart function counts in is not fixed by the text. We allow every lattice
$\Lambda\le X$ that contains $\lambda$, which is exactly the condition that $P_\lambda$ is a lattice
polytope for $\Lambda$. This includes the full (co)weight lattice $\Lambda = X$, the root lattice
$\Lambda = \mathbb{Z}\Phi$ when $\lambda\in\mathbb{Z}\Phi$, and the lattice $\mathbb{Z}\lambda$
generated by the vertices.
(ii) If one counts in the coset $t\lambda + \mathbb{Z}\Phi$ (the weights of $V(t\lambda)$), this is
the root-lattice count whenever $\lambda\in\mathbb{Z}\Phi$, so it is covered for such $\lambda$, for
instance $\lambda = a$, the root. We do \emph{not} cover this coset count for
$\lambda\notin\mathbb{Z}\Phi$. Indeed, for the $SL_2$ datum and $\lambda = 1$ it gives $t+1$, which
\emph{is} monic of degree $1$. The law is still refuted under this reading, by $\lambda = a$.
(iii) A polynomial ``is'' the Ehrhart function if it agrees with it at every positive integer $t$. We
allow any real polynomial, which includes the rational ones. The hypothesis is only required for
$t\ge 1$, so the statement is not weakened by the convention $L(P,0) = 1$.
(iv) We take $\lambda\ne 0$. For $\lambda = 0$ the polytope is a point and the case is degenerate.

\section{The theorem}

\begin{theorem}\label{thm:main}
Let $ab = 2$, let $0\ne\lambda\in X = \mathbb{Z}$, and let $\Lambda\le\mathbb{Z}$ be a lattice with
$\lambda\in\Lambda$. For $\mathcal{R}(a,b)$:
\begin{enumerate}
\item $W$ acts on $X$ as $\{\pm 1\}$, so $W\cdot\lambda = \{\lambda,-\lambda\}$ and
$P_\lambda = [-|\lambda|,|\lambda|]$.
\item $L_X(P_\lambda,t) = 2|\lambda|\,t + 1$ for all $t$. More generally, if $\Lambda = d\mathbb{Z}$
then $L_\Lambda(P_\lambda,t) = 2c\,t+1$ with $c = |\lambda/d|\ge 1$.
\item No real polynomial $p$ with $p(t) = L_\Lambda(P_\lambda,t)$ for all integers $t\ge 1$ is monic,
of any degree. In particular, none is monic of degree $1$, the semisimple rank.
\end{enumerate}
Hence the semisimple rank law fails for $SL_2$ and for $PGL_2$, for every nonzero $\lambda$ and every
lattice convention in (i). So the conjecture, being a conjunction, is false.
\end{theorem}

\begin{proof}
(1) For either root $\pm a$ with coroot $\pm b$ we have
$s(x) = x - (x\cdot(\pm b))(\pm a) = x - abx = -x$. Every element of $W$ is a product of reflections,
so by induction on the generation of $W$ it acts as $+1$ or $-1$. The reflection lies in $W$, so
$W\cdot\lambda = \{\lambda,-\lambda\}$. The convex hull of two points of $\mathbb{R}$ is the segment
between them, which is $[-|\lambda|,|\lambda|]$.

(2) For $t\ge 0$ we have $tP_\lambda = [-t|\lambda|, t|\lambda|]$. Every subgroup of $\mathbb{Z}$ is
$d\mathbb{Z}$ for some $d$ (the ideal $\Lambda$ of the PID $\mathbb{Z}$ is principal), and
$\lambda\in\Lambda$ gives $\lambda = dm$ with $d,m\ne0$. The map $k\mapsto dk$ is a bijection from
$\{k : |k|\le t|m|\}$ onto $\{x\in d\mathbb{Z} : |x|\le t|dm|\}$, so
$L_\Lambda(P_\lambda,t) = 2|m|t+1$. For $d = 1$ this is $2|\lambda|t+1$.

(3) If $p(t) = 2ct+1$ for infinitely many $t$, then $p = 2cX + 1$, because two polynomials that agree
at infinitely many points are equal. Its leading coefficient is $2c\ge 2$, which is not $1$.
\end{proof}

For example, for $SL_2$ and $\lambda = 1$ (the fundamental weight) in the weight lattice we get
$L(t) = 2t+1$. For $\lambda = a$ (the root) in the root lattice we get $L(t) = 2t+1$ for both data.
Two values already rule out a monic linear $t+c$, since it would force $L(2)-L(1) = 1$.

\begin{remark}[Second law; not formalized, not used]\label{rem:clause2}
For the $SL_2$ datum, $\operatorname{vol}(P_n) = 2n$, while $\dim V(n)/|W| = (n+1)/2$. The ratio
$4n/(n+1)$ is not constant in $n$, so no fixed normalization makes the second law true either. This
remark is outside the Lean development, and the disproof does not depend on it.
\end{remark}

\begin{remark}[Scope]
Only rank one is formalized. A computer check (in the attempt record, not in Lean) for $A_2$ with
$\lambda = \omega_1$ in the weight lattice gives $L(1),\dots,L(4) = 4,10,19,31$, whose leading
coefficient is $3/2$.
\end{remark}

\section{Lean formalization}

The file \texttt{lean/Conjecture8228/Basic.lean} (Lean 4, Mathlib \texttt{v4.33.1}, 268 lines) is built
with \texttt{lake build}. Its main declarations are:
\begin{itemize}
\item \texttt{C8228.rk1 a b hab : RootDatum (Fin 2) $\mathbb{Z}$ $\mathbb{Z}$}, with
\texttt{rootDatumSL2 = rk1 2 1}, \texttt{rootDatumPGL2 = rk1 1 2}. The perfect pairing is
\texttt{LinearMap.mul}.
\item \texttt{weyl\_smul}, \texttt{orbit\_eq}: $W$ acts as $\pm1$, and the orbit is $\{\lambda,-\lambda\}$.
\item \texttt{weightPolytope\_eq}: $P_\lambda = $ \texttt{Icc (-|lam|) |lam|}.
\item \texttt{semisimpleRank\_rk1}: the semisimple rank is $1$.
\item \texttt{ehrhart\_top}: $L_X(P_\lambda,t) = 2|\lambda|t+1$ for every $\lambda$ and $t$.
\item \texttt{ehrhart\_formula}: for $\lambda\ne0$ and any lattice $\Lambda\ni\lambda$,
$L_\Lambda(P_\lambda,t) = 2ct+1$ with $c\ge1$.
\item \texttt{ehrhart\_not\_monic} (main): for \texttt{lam $\ne$ 0} and \texttt{lam $\in$ Lam},\\
\texttt{$\neg\exists$ p : $\mathbb{R}$[X], p.Monic $\wedge$ $\forall$ t : $\mathbb{N}$, 0 < t $\to$ p.eval t = ehrhart Lam (weightPolytope (rk1 a b hab) lam) t}.
\item \texttt{semisimple\_rank\_law\_fails}: the same statement with the extra conjunct
\texttt{p.natDegree = semisimpleRank (rk1 a b hab)}. Its specialization to the root lattice
\texttt{(rk1 a b hab).rootSpan $\mathbb{Z}$} is \texttt{semisimple\_rank\_law\_fails\_rootLattice}.
\end{itemize}

\paragraph{Axiom audit.}
\texttt{\#print axioms} reports \texttt{[propext, Classical.choice, Quot.sound]} for
\texttt{ehrhart\_not\_monic}, \texttt{semisimple\_rank\_law\_fails},
\texttt{semisimple\_rank\_law\_fails\_rootLattice}, \texttt{ehrhart\_top} and
\texttt{semisimpleRank\_rk1}. There is no \texttt{sorry} and no \texttt{native\_decide}, and no axioms
are added.

\begin{thebibliography}{9}
\bibitem{wiki} Wikipedia, \emph{Ehrhart polynomial},
\url{https://en.wikipedia.org/wiki/Ehrhart_polynomial} (retrieved 2026-10-05).
\end{thebibliography}

\end{document}
